\documentclass[12pt,a4paper]{article}
\usepackage[utf8]{inputenc}
\usepackage[T1]{fontenc}
\usepackage{amsmath, amsfonts, amssymb, amsthm}
\usepackage{graphicx}
\usepackage{listings}
\usepackage{xcolor}
\usepackage{hyperref}
\usepackage{geometry}
\usepackage{fancyhdr}
\usepackage{titlesec}
\usepackage{tocloft}

\geometry{margin=1in}

% --- Custom Colors ---
\definecolor{codegreen}{rgb}{0,0.6,0}
\definecolor{codegray}{rgb}{0.5,0.5,0.5}
\definecolor{codepurple}{rgb}{0.58,0,0.82}
\definecolor{backcolour}{rgb}{0.95,0.95,0.92}
\definecolor{navy}{rgb}{0,0,0.5}

% --- Header and Footer ---
\pagestyle{fancy}
\fancyhf{}
\rhead{CollabVerse: Technical Report}
\lhead{DAA \& CP Lab}
\cfoot{\thepage}
\setlength{\headheight}{14.5pt}

% --- Listings Style ---
\lstset{
    backgroundcolor=\color{backcolour},   
    commentstyle=\color{codegreen},
    keywordstyle=\color{magenta},
    numberstyle=\tiny\color{codegray},
    stringstyle=\color{codepurple},
    basicstyle=\ttfamily\footnotesize,
    breakatwhitespace=false,         
    breaklines=true,                 
    captionpos=b,                    
    keepspaces=true,                 
    numbers=left,                    
    numbersep=5pt,                  
    showspaces=false,                
    showstringspaces=false,
    showtabs=false,                  
    tabsize=2
}

\title{
    \vspace{2cm}
    \textbf{BACHELORS OF TECHNOLOGY}\\
    \large{COMPUTER SCIENCE AND ENGINEERING}\\
    \vspace{1cm}
    \huge{\textbf{CollabVerse: Intelligent Network Connector}}\\
    \vspace{0.5cm}
    \Large{A Comprehensive Technical Report on Design and Analysis of Algorithms}\\
    \vspace{2cm}
}

\author{
    \textbf{Submitted By:}\\
    \large{Sansriti Mishra (2401030293)}\\
    \large{Sanvi Dhingra (2401030310)}\\
    \large{Anant Kumar Verma (2401030295)}\\
    \vspace{1cm}\\
    \textbf{Submitted To:}\\
    \large{Dr. Kirti Jain}\\
    \large{Dr. Aastha Maheswari}
}

\date{Academic Year 2025-2026}

\begin{document}

\maketitle
\newpage

\begin{abstract}
CollabVerse is an advanced team formation and collaboration platform designed to optimize resource allocation within academic and professional communities. The system employs a hybrid architecture where a FastAPI (Python) backend handles web requests, authentication, and database operations, while a high-performance C++ engine executes computationally intensive algorithms including Dinic's Algorithm for maximum bipartite matching, Bitmask Dynamic Programming for synergy optimization, and a Multi-Layer Perceptron (MLP) neural network for skill prediction from GitHub metadata. This report provides an in-depth exploration of the system's architecture, the mathematical foundations of its core algorithms, and a rigorous analysis of their time and space complexities. The implementation demonstrates how theoretical algorithmic concepts from Design and Analysis of Algorithms (DAA) and Competitive Programming (CP) can be bridged with modern web technologies to create a scalable, real-time collaboration ecosystem.
\end{abstract}

\newpage
\tableofcontents
\newpage
\listoffigures
\listoftables
\newpage

\section{Introduction}

\subsection{Overview}
CollabVerse is an intelligent team formation and collaboration platform designed to address the longstanding challenge of optimal resource allocation within academic and professional communities. Unlike conventional networking platforms that rely on self-declared skills which are often inaccurate or outdated, CollabVerse extracts authentic technical proficiency directly from GitHub metadata, including repository languages, commit frequency, and project complexity. The system leverages advanced graph algorithms from the Design and Analysis of Algorithms (DAA) domain, specifically implementing Dinic's Algorithm for maximum bipartite matching to solve the team allocation problem optimally. Additionally, it employs Bitmask Dynamic Programming for synergy optimization when forming teams for individual projects with a limited number of roles. The platform is built on a modern technology stack comprising FastAPI (Python) for backend services, SQLite for lightweight data persistence, and React with Tailwind CSS for an interactive, visually engaging frontend.

\subsection{Motivation}
The primary motivation behind CollabVerse is to eliminate manual bias and social influence from the team formation process. In typical academic environments, students with similar skill sets tend to cluster together, while interdisciplinary projects that require diverse expertise struggle to find suitable connections. By applying rigorous graph theory and network flow algorithms, CollabVerse ensures that every project role receives the most compatible candidate, skill diversity is maximized within teams, and the overall allocation is globally optimal rather than locally greedy. Furthermore, by grounding skill assessment in actual GitHub activity rather than self-reported surveys, the platform provides an objective, verifiable basis for compatibility scoring.

\subsection{Project Scope}
The CollabVerse project encompasses four major functional components. First, the authentication and verification module handles secure user registration, login with JWT tokens, and a novel GitHub-based proof-of-ownership verification system. Second, the skill ingestion module recursively fetches GitHub repository data to build comprehensive user skill profiles based on actual coding activity. Third, the team formation engine implements Dinic's Max Flow algorithm in C++ to solve the student-to-project assignment problem optimally. Fourth, the collaboration layer provides real-time WebSocket-based chat, team dashboards, and data-driven insights. Together, these components form a complete, production-ready platform for intelligent team formation.

\newpage
\section{Literature Review \& Algorithmic Background}

\subsection{Network Flow Theory}
Network flow theory represents a cornerstone of combinatorial optimization and has found extensive applications in matching, scheduling, and resource allocation problems. In a directed graph \( G = (V, E) \) with a designated source node \( s \) and sink node \( t \), each edge \( (u,v) \) is assigned a non-negative capacity \( c(u,v) \) representing the maximum amount of flow that can pass through that edge. The maximum flow problem seeks to determine the greatest possible amount of flow that can be sent from \( s \) to \( t \) without violating any edge capacity constraints. This formulation elegantly captures many real-world allocation problems, including the team formation challenge addressed by CollabVerse.

\subsubsection{Ford-Fulkerson Method}
The Ford-Fulkerson method represents the earliest approach to solving maximum flow problems. The algorithm iteratively augments flow along paths found in the residual graph, which records remaining capacities and reverse flow opportunities. While conceptually simple, the Ford-Fulkerson method suffers from a time complexity of \( O(E \cdot f) \), where \( f \) is the maximum flow value. This complexity is pseudo-polynomial because \( f \) can be arbitrarily large, potentially leading to extremely slow convergence when capacities are large integers.

\subsubsection{Edmonds-Karp Algorithm}
The Edmonds-Karp algorithm improves upon Ford-Fulkerson by always selecting the shortest augmenting path, measured by the number of edges, using Breadth-First Search (BFS). This seemingly minor modification yields a polynomial time complexity of \( O(V E^2) \), which is a significant theoretical improvement. However, for dense graphs with thousands of vertices, this complexity remains impractical for real-time applications.

\subsubsection{Dinic's Algorithm}
Dinic's algorithm represents a substantial advancement in maximum flow computation. The algorithm introduces the concept of a level graph, also known as a layered network, constructed via BFS from the source. Each node is assigned a level representing its shortest distance from the source. Subsequently, multiple Depth-First Search (DFS) passes are executed to find blocking flows within the level graph. The key insight is that each blocking flow increases the level of the sink, ensuring that at most \( V \) phases are required. The overall complexity of Dinic's algorithm is \( O(E V^2) \) for general graphs. More importantly for CollabVerse, when applied to bipartite matching networks, the complexity reduces to \( O(E \sqrt{V}) \), making it the algorithm of choice for our team formation problem.

\subsection{Bipartite Matching as Maximum Flow}
The student-to-project assignment problem naturally maps to a bipartite matching formulation. The left partition consists of students \( S = \{s_1, s_2, \ldots, s_n\} \), while the right partition comprises project roles \( R = \{r_1, r_2, \ldots, r_m\} \). An edge exists between student \( s_i \) and role \( r_j \) if the student possesses the required skills for that role, as determined by the GitHub skill ingestion module. By adding a super-source node connected to every student with capacity one, and a super-sink node connected from every role with capacity one, the maximum flow from source to sink equals the size of the maximum matching. This transformation allows us to leverage Dinic's algorithm to solve the team formation problem with provable optimality guarantees.

\subsection{Combinatorial Optimization with Bitmask Dynamic Programming}
When forming a team for a single project with a limited number of roles, typically fifteen or fewer, we can employ Bitmask Dynamic Programming to achieve optimal synergy maximization. The state is represented by a bitmask where each bit indicates whether a particular role has been filled. The DP recurrence considers adding students one by one, either skipping a student or assigning them to an empty role. The time complexity of this approach is \( O(N \cdot 2^R) \), which, while exponential in the number of roles, remains practical for the small role sets typical of academic projects. This method is particularly valuable when project leaders want to form the best possible team from a pool of candidates rather than solving the global allocation problem.

\subsection{Neural Networks for Skill Prediction}
CollabVerse employs a Multi-Layer Perceptron (MLP) neural network implemented in C++ to predict a developer's expertise level across various programming languages and frameworks. The network takes as input a high-dimensional feature vector derived from GitHub metadata, including repository counts per language, average commit frequency, repository stars and forks, and contribution recency. The MLP consists of an input layer matching the feature dimension, two hidden layers using sigmoid activation functions, and an output layer that produces a probability distribution over skill proficiency levels. The network is trained offline using a labeled dataset of GitHub profiles and then integrated into the C++ engine for real-time inference.

\subsection{Greedy Paradigms for Real-Time Matching}
In scenarios where instantaneous response is prioritized over global optimality, CollabVerse utilizes greedy algorithmic paradigms. A greedy algorithm makes the locally optimal choice at each stage with the hope of finding a global optimum. In our context, the "Build for Me" feature implements a greedy selection process that ranks potential teammates by a complementarity score and selects the top candidates immediately. This ensures that the user receives an interactive experience without the latency associated with global graph re-computation.

\subsection{Algorithmic Paradigms Summary}
The CollabVerse algorithmic engine is a hybrid of three primary paradigms:
\begin{itemize}
    \item \textbf{Optimization (Graph Theory)}: Dinic's Algorithm is used to solve the Maximum Bipartite Matching problem, ensuring a globally optimal assignment of students to projects.
    \item \textbf{Dynamic Programming (Bitmasking)}: Used for synergy maximization in small teams, reducing an $O(N!)$ search space to $O(N \cdot 2^R)$ via memoization.
    \item \textbf{Greedy (Heuristics)}: Used in the Python-based recommendation service to provide sub-second teammate suggestions based on local complementarity scores.
\end{itemize}

\subsection{Advanced UI Analytics}
To help users understand the collaboration landscape, CollabVerse employs a Bento-style dashboard with sophisticated data visualizations. The interface provides real-time metrics on team strength, complementarity scores, and skill coverage. These visualizations are built using modern SVG-based gauge components and glassmorphic skill bars, enabling users to identify critical gaps and synergy opportunities at a glance.

\newpage
\section{Problem Formulation}

\subsection{The Global Team Allocation Problem}
The global team allocation problem can be formally stated as follows. Given a set of students \( S = \{s_1, s_2, \ldots, s_n\} \) and a set of projects \( P = \{p_1, p_2, \ldots, p_k\} \), where each project \( p_j \) requires a specific set of roles \( R_{p_j} = \{r_1, r_2, \ldots, r_{m_j}\} \), we define a compatibility score \( \text{comp}(s_i, r_j) \in [0,1] \) derived from the student's GitHub activity. The objective is to find an assignment \( A: S \rightarrow (P \times R) \) that maximizes the sum of compatibility scores \( \sum \text{comp}(s_i, \text{assigned role}) \), subject to three constraints: each student may be assigned to at most one role, each role may be assigned to at most one student, and all required roles for a project should be filled if a sufficient number of qualified students exist. This formulation captures the essence of fair and optimal team formation.

\subsection{The Local Team Formation Problem}
For scenarios where a project leader wishes to form a single team from a pool of candidates, we consider a slightly different optimization objective. Given a specific project requiring a set of roles, we seek to find a subset of students \( S' \subseteq S \) that maximizes a composite score function: \( \text{Score}(S') = \sum_{s \in S'} \text{comp}(s, \text{role}(s)) + \lambda \cdot \text{diversity}(S') \), where diversity is measured as the number of unique programming languages and frameworks represented in the team. The parameter \( \lambda \) allows the project leader to tune the trade-off between individual competence and team diversity. This problem is solved using Bitmask Dynamic Programming when the number of roles is small, or through heuristic approaches when the role count exceeds the feasible limit for exact DP.

\newpage
\section{System Architecture}

\subsection{High-Level Architecture}
CollabVerse follows a hybrid modular architecture that separates the web-facing components from the computationally intensive algorithm engine. The React-based frontend communicates with the FastAPI backend via HTTP REST endpoints for data operations and WebSocket connections for real-time messaging. The backend handles authentication, business logic, and database interactions using SQLite as the persistent storage layer. When algorithm execution is required, the Python backend invokes a compiled C++ binary that performs the core optimization tasks including maximum flow computation, Bitmask DP, and neural network inference. This hybrid approach combines the rapid development and rich ecosystem of Python with the raw performance of C++ for compute-intensive operations.

\subsection{Backend Framework: FastAPI}
FastAPI was selected as the backend framework for several compelling reasons. First, its native asynchronous support is essential for non-blocking GitHub API calls and concurrent WebSocket connection management. Second, FastAPI automatically generates OpenAPI documentation from Python type hints, which significantly simplifies frontend-backend integration and API testing. Third, independent benchmarks demonstrate that FastAPI achieves performance comparable to Node.js and Go, making it suitable for production deployment. The backend is organized into several modules: the authentication module handles JWT token generation and bcrypt password hashing; the WebSocket manager maintains thread-safe registries of active connections; the GitHub client provides asynchronous methods for repository fetching; and the C++ bridge module manages subprocess execution of the performance engine.

\subsection{C++ Performance Engine}
The core algorithmic components of CollabVerse are implemented in C++ for maximum performance. This engine is compiled into an optimized binary that is invoked by the Python backend whenever team formation or skill prediction is required. The engine implements Dinic's Algorithm for maximum flow, Bitmask Dynamic Programming for synergy optimization, and an MLP neural network for skill prediction. Communication between Python and C++ occurs via standard input/output streams using a simple protocol: the Python backend serializes input data as JSON, passes it to the C++ binary, and parses the JSON output containing the computed results. This architecture ensures that computationally expensive operations do not block the asynchronous web server.

\subsection{Frontend Framework: React and Tailwind CSS}
The frontend is built using React, a component-based library for building user interfaces. The dashboard features several key visualizations, including Bento-grid layouts that render the user's technical profile. Profile cards are color-coded by primary proficiency, while progress bars represent skill development levels. Tailwind CSS is used throughout to implement a glassmorphic design system characterized by frosted glass panels, subtle gradients, and neon accents. React Router manages navigation between different views including user profiles, team dashboards, and chat interfaces. Framer Motion provides smooth micro-interactions that enhance user engagement without sacrificing performance.

\subsection{Database Design}
CollabVerse uses SQLite for its simplicity, zero-configuration deployment, and sufficient performance for the expected scale of academic use. The database schema consists of several interrelated tables. The users table stores authentication credentials and profile information including username, email address, and hashed password. The github\_verification table tracks the verification codes issued to users along with timestamps to enable proof-of-ownership validation. The skills table maintains a many-to-many relationship between users and programming languages, with an associated experience level derived from repository analysis. The teams table stores metadata about formed teams including project name and creation timestamp. The team\_members table maps users to specific teams and records the role each member has been assigned. Finally, the chat\_messages table provides persistent storage for all team communication, enabling users to revisit conversations even after disconnecting.

\newpage
\section{Core Algorithmic Implementation}

\subsection{Dinic's Algorithm in C++}
The implementation of Dinic's algorithm in CollabVerse's C++ engine follows the classical approach with careful attention to data structure efficiency. The implementation uses adjacency lists with edge structures that store destination vertex, capacity, flow, and the index of the reverse edge for residual graph updates. The BFS phase constructs level graphs using a queue-based approach, while the DFS phase uses a pointer array to avoid re-examining saturated edges. The implementation is optimized for bipartite matching problems where the graph structure is known to be bipartite, allowing for the \( O(E \sqrt{V}) \) complexity guarantee.

\begin{lstlisting}[language=C++, caption=Dinic's Algorithm Implementation in C++]
struct Edge {
    int to;
    long long cap, flow;
    int rev;
};

class DinicSolver {
    vector<vector<Edge>> adj;
    vector<int> level, ptr;
    
    bool bfs(int s, int t) {
        fill(level.begin(), level.end(), -1);
        level[s] = 0;
        queue<int> q;
        q.push(s);
        while (!q.empty()) {
            int v = q.front(); q.pop();
            for (auto& edge : adj[v]) {
                if (edge.cap - edge.flow > 0 && level[edge.to] == -1) {
                    level[edge.to] = level[v] + 1;
                    q.push(edge.to);
                }
            }
        }
        return level[t] != -1;
    }
    
    long long dfs(int v, int t, long long pushed) {
        if (pushed == 0 || v == t) return pushed;
        for (int& cid = ptr[v]; cid < (int)adj[v].size(); ++cid) {
            auto& edge = adj[v][cid];
            int tr = edge.to;
            if (level[v] + 1 != level[tr] || edge.cap - edge.flow == 0) continue;
            long long tr_pushed = dfs(tr, t, min(pushed, edge.cap - edge.flow));
            if (tr_pushed == 0) continue;
            edge.flow += tr_pushed;
            adj[tr][edge.rev].flow -= tr_pushed;
            return tr_pushed;
        }
        return 0;
    }
    
public:
    DinicSolver(int n) : adj(n), level(n), ptr(n) {}
    
    void add_edge(int from, int to, long long cap) {
        adj[from].push_back({to, cap, 0, (int)adj[to].size()});
        adj[to].push_back({from, 0, 0, (int)adj[from].size() - 1});
    }
    
    long long max_flow(int s, int t) {
        long long flow = 0;
        while (bfs(s, t)) {
            fill(ptr.begin(), ptr.end(), 0);
            while (long long pushed = dfs(s, t, 1e18)) {
                flow += pushed;
            }
        }
        return flow;
    }
};
\end{lstlisting}

\subsection{Bitmask Dynamic Programming for Synergy}
The synergy optimization engine in C++ uses Bitmask DP to find the optimal assignment of students to project roles. The implementation uses memoization with a two-dimensional DP table where the first dimension indexes the student being considered and the second dimension is a bitmask representing filled roles. The recurrence relation chooses between skipping the current student or assigning them to any unfilled role. The base case occurs when all roles are filled or no students remain.

\begin{lstlisting}[language=C++, caption=Bitmask DP Synergy Optimization in C++]
class SynergyOptimizer {
    int num_students, num_roles;
    vector<vector<int>> skill_matrix;
    vector<vector<int>> memo;
    
    int solve(int s_idx, int mask) {
        if (mask == (1 << num_roles) - 1) return 0;
        if (s_idx == num_students) return -1e9;
        if (memo[s_idx][mask] != -1) return memo[s_idx][mask];
        
        int res = solve(s_idx + 1, mask);
        
        for (int r = 0; r < num_roles; ++r) {
            if (!(mask & (1 << r))) {
                res = max(res, skill_matrix[s_idx][r] + 
                         solve(s_idx + 1, mask | (1 << r)));
            }
        }
        return memo[s_idx][mask] = res;
    }
    
public:
    SynergyOptimizer(int s, int r, const vector<vector<int>>& matrix)
        : num_students(s), num_roles(r), skill_matrix(matrix) {
        memo.assign(s, vector<int>(1 << r, -1));
    }
    
    int getBestSynergy() {
        return solve(0, 0);
    }
};
\end{lstlisting}

\subsection{Neural Network MLP Predictor}
The C++ engine implements a feedforward neural network for skill prediction from GitHub features. The MLP architecture consists of an input layer matching the feature vector dimension, multiple hidden layers using sigmoid activation, and an output layer producing skill proficiency probabilities. The implementation uses standard C++ vectors for weight matrices and biases, providing a lightweight yet powerful inference engine.

\begin{lstlisting}[language=C++, caption=MLP Neural Network Implementation in C++]
class NeuralNetwork {
    struct Layer {
        vector<vector<float>> weights;
        vector<float> biases;
        vector<float> outputs;
        vector<float> deltas;
        
        Layer(int input_size, int output_size);
    };
    
    vector<Layer> layers;
    float learning_rate;
    
    float sigmoid(float x) {
        return 1.0f / (1.0f + std::exp(-x));
    }
    
public:
    NeuralNetwork(const vector<int>& layer_sizes, float lr = 0.01);
    
    vector<float> forward(const vector<float>& input) {
        vector<float> current = input;
        for (auto& layer : layers) {
            vector<float> next(layer.weights.size(), 0.0f);
            for (size_t i = 0; i < layer.weights.size(); i++) {
                float sum = layer.biases[i];
                for (size_t j = 0; j < current.size(); j++) {
                    sum += layer.weights[i][j] * current[j];
                }
                next[i] = sigmoid(sum);
            }
            layer.outputs = next;
            current = next;
        }
        return current;
    }
    
    vector<float> predict(const vector<float>& input) {
        return forward(input);
    }
};
\end{lstlisting}

\subsection{GitHub Skill Ingestion in Python}
The skill ingestion module is implemented in Python and is responsible for converting a GitHub username into a structured skill profile. The process begins by calling the GitHub REST API endpoint that lists all public repositories belonging to the user. The API response is paginated, so the implementation handles pagination by checking for the \texttt{Link} header and following \texttt{next} pages until all repositories have been fetched. For each repository, the language field is extracted and counted. After processing all repositories, the raw counts are normalized into feature vectors that can be passed to the C++ MLP engine for skill prediction.

\begin{lstlisting}[language=Python, caption=GitHub Skill Ingestion in Python]
async def fetch_user_skills(username: str):
    url = f"https://api.github.com/users/{username}/repos?per_page=100"
    skills = {}
    
    async with httpx.AsyncClient() as client:
        response = await client.get(url, auth=(GITHUB_TOKEN, ""))
        repos = response.json()
        
        for repo in repos:
            lang = repo.get("language")
            if lang:
                skills[lang] = skills.get(lang, 0) + 1
    
    total = sum(skills.values())
    if total == 0:
        return {}
    return {lang: count/total for lang, count in skills.items()}
\end{lstlisting}

\newpage
\section{The GitHub Verification Flow}

\subsection{Problem Statement}
One of the fundamental challenges in building CollabVerse was establishing a trustworthy identity verification mechanism. In an open platform, any user could claim any GitHub username without providing proof of ownership. Traditional solutions such as OAuth authentication, while secure, require users to grant extensive permissions to the application, which can be a barrier to adoption. The verification system needed to be decentralized, cryptographically secure, and minimally invasive to user privacy.

\subsection{The Proof of Repository Challenge}
CollabVerse implements a novel verification scheme based on a proof-of-repository challenge. When a user initiates verification for a claimed GitHub username, the backend generates a random eight-character alphanumeric code and stores it in the database with an associated expiry timestamp. The user is then instructed to create a public repository named exactly \texttt{CollabVerse-Verification} and add a file named \texttt{verify.txt} containing the generated code.

After the user completes these steps and clicks the verify button, the backend fetches the file content using the GitHub REST API. The API returns the file content in Base64-encoded format, which the backend decodes and compares against the stored code. If the codes match, the user's account is marked as verified, and the verification timestamp is recorded. This approach ensures that only the legitimate owner of the GitHub account can complete verification, as only they can create repositories under that account.

\subsection{Bypassing CDN Cache}
An important technical challenge encountered during implementation was GitHub's aggressive content delivery network (CDN) caching. When a file is accessed via the raw.githubusercontent.com domain, the response may be cached for several minutes, causing verification to fail even after the user has correctly created the file. To solve this, the verification endpoint was modified to use the GitHub REST API with specific HTTP headers that force a fresh fetch from the master branch. The \texttt{Accept: application/vnd.github.v3.raw} header instructs the API to return the raw file content, while the \texttt{If-None-Match: ""} header overrides conditional caching. This combination guarantees that the backend receives the most current version of the verification file.

\begin{lstlisting}[language=Python, caption=GitHub Verification with Cache Bypass]
headers = {
    "Accept": "application/vnd.github.v3.raw",
    "If-None-Match": ""
}
url = f"https://api.github.com/repos/{username}/CollabVerse-Verification/contents/verify.txt"
response = await client.get(url, headers=headers)
\end{lstlisting}

\newpage
\section{Complexity Analysis}

\subsection{Time Complexity}
The time complexity of CollabVerse varies across its different modules, each optimized for its specific problem domain. The global team matching module implements Dinic's algorithm in C++ for maximum flow on a bipartite graph. For a bipartite graph with \(V\) vertices and \(E\) edges, Dinic's algorithm achieves a complexity of \(O(E \sqrt{V})\), which is significantly better than the \(O(V E^2)\) complexity of Edmonds-Karp for dense graphs. This improvement is achieved through the level graph construction and the use of blocking flows.

The local team formation module employs Bitmask Dynamic Programming in C++ with a complexity of \(O(N \cdot 2^R)\), where \(N\) is the number of candidate students and \(R\) is the number of required roles. While this complexity is exponential in the number of roles, the constant factor is small, and the approach remains practical for \(R \leq 15\), which covers the vast majority of academic projects.

The neural network MLP predictor performs inference in \(O(\sum_{i=1}^{L} (n_{i-1} \cdot n_i))\) time, where \(L\) is the number of layers and \(n_i\) is the number of neurons in layer \(i\). With the current architecture of input layer dimension 50, two hidden layers of 128 and 64 neurons, and an output layer of 10 skill levels, each inference requires approximately 15,000 floating-point operations, executing in microseconds. The skill ingestion module makes sequential calls to the GitHub API, with the number of calls proportional to the number of repositories belonging to the user. For users with typical activity, this completes within a few seconds.

\begin{table}[h!]
\centering
\begin{tabular}{|l|l|l|}
\hline
\textbf{Module} & \textbf{Implementation} & \textbf{Complexity} \\ \hline
Team Matching (Global) & C++ (Dinic) & \(O(E \sqrt{V})\) \\ \hline
Team Formation (Local) & C++ (Bitmask DP) & \(O(N \cdot 2^R)\) \\ \hline
Greedy Team Matching & Python (Heuristic) & \(O(N \log N)\) \\ \hline
Skill Prediction & C++ (MLP) & \(O(\sum n_{i-1} n_i)\) \\ \hline
Skill Ingestion & Python (GitHub API) & \(O(\text{Repositories})\) \\ \hline
Frontend Rendering & React & \(O(N)\) \\ \hline
Password Hashing & Python (Bcrypt) & \(O(2^k)\) \\ \hline
Chat Messaging & Python (WebSocket) & \(O(\text{team\_size})\) \\ \hline
\end{tabular}
\caption{System Time Complexity Analysis}
\label{tab:complexity}
\end{table}

\subsection{Space Complexity}
The space complexity of CollabVerse is dominated by several data structures. The flow network for global matching is stored as an adjacency list requiring \(O(V + E)\) space, which is efficient for the sparse graphs typical of academic skill networks. The Bitmask DP table for local synergy optimization requires \(O(N \cdot 2^R)\) space, which can become substantial when both \(N\) and \(R\) are large. In practice, the implementation restricts \(R\) to fifteen or fewer, keeping memory usage manageable.

The MLP neural network stores weight matrices of total size \( \sum_{i=1}^{L} (n_{i-1} \cdot n_i) \), approximately 15,000 parameters in the current configuration, which occupies negligible memory. The WebSocket connection manager maintains references to all active connections, requiring \(O(U)\) space where \(U\) is the number of concurrently online users. The React frontend stores the local state for dashboards and chat, requiring \(O(V + E)\) space for normalized data. The SQLite database file size scales linearly with the number of users, teams, and messages stored.

\subsection{Optimization Techniques}
Several optimization techniques were employed to achieve the stated complexities. The level graph construction in Dinic's algorithm prunes the search space by restricting flow to monotonically increasing levels, preventing cycles and redundant exploration. The current edge pointer array in C++ prevents re-examination of edges that have already been saturated in the current phase, which is essential for achieving the optimal theoretical complexity. GitHub data is cached in the database with a 24-hour time-to-live to reduce API calls and improve response times. Multiple users' GitHub data can be fetched concurrently using Python's \texttt{asyncio.gather}, reducing total latency when processing multiple skill ingestion requests. The C++ engine is compiled with \texttt{-O3} optimization flags to ensure maximum execution speed.

\newpage
\section{Testing \& Quality Assurance}

\subsection{Unit Testing}
The CollabVerse codebase includes comprehensive unit tests for all algorithmic components. The C++ Dinic implementation is tested on standard bipartite matching problems drawn from competitive programming platforms such as Codeforces and CSES. These test cases include perfect matching scenarios where all roles can be filled, imperfect matching scenarios where some roles remain unfilled, and edge cases involving zero-capacity edges or disconnected components. Each test asserts that the maximum flow computed by the algorithm matches the expected value derived from independent verification.

\subsection{Integration Testing}
Beyond unit tests, CollabVerse undergoes rigorous integration testing to validate end-to-end workflows. The complete user journey is tested from registration through team formation and chat. This includes user registration followed by login and JWT issuance; GitHub verification followed by skill extraction and profile completion; team creation followed by C++ engine invocation and member assignment; and WebSocket connection followed by message broadcast and delivery confirmation. These integration tests run continuously to catch regressions introduced by code changes.

\subsection{Edge Case Analysis}
Special attention was paid to edge cases that could cause system failures or poor user experiences. When a user has no public GitHub repositories, the system gracefully falls back to manually entered skills collected during registration. When there are more students than available roles, Dinic's algorithm naturally produces a matching where some students remain unassigned, and the system notifies these students while suggesting alternative projects. When there are more roles than students, some project roles remain unfilled, and the system alerts the project creator to adjust requirements or recruit additional members.

Concurrency issues are handled through appropriate locking mechanisms. The WebSocket manager uses asyncio locks to protect shared data structures during concurrent connection establishment and message broadcasting. The SQLite database is configured with the appropriate pragmas to handle concurrent reads and writes from multiple FastAPI worker processes. The C++ engine is invoked as a subprocess for each request, ensuring isolation between concurrent team formation requests. GitHub API rate limits are managed through exponential backoff and caching, ensuring that the system remains functional even under heavy usage.

\newpage
\section{UI/UX Design Philosophy}

\subsection{Visual Design Language}
The visual design of CollabVerse follows a glassmorphic aesthetic that has become increasingly popular in modern web applications. Glassmorphism is characterized by frosted glass panels with background blur, subtle gradients, and neon accent colors. This design language was chosen because it creates a premium, modern feel while maintaining excellent readability. The dark mode first approach recognizes that developers typically prefer low-light environments and spend long hours working in code editors with dark themes.

\subsection{Bento-Grid Analytics}
The centerpiece of the CollabVerse dashboard is the Bento-grid analytics system. This layout organizes complex data into a clean, modular interface where each "bento box" represents a specific metric or insight. Users can quickly scan their technical proficiency, team complementarity scores, and upcoming project milestones. This approach maximizes information density while minimizing cognitive load, providing a streamlined experience for team leaders and developers alike.

\subsection{Responsive Design}
The CollabVerse interface is fully responsive, adapting gracefully to different screen sizes and input methods. On mobile devices, a collapsible sidebar conserves screen real estate while maintaining access to navigation. On tablets, the interface switches to a two-column layout that balances information density with readability. On desktop systems with large displays, the dashboard expands to a multi-panel layout with fixed navigation and a spacious activity feed. Tailwind CSS utility classes provide the foundation for this responsive behavior, with breakpoints at standard device widths.

\newpage
\section{Future Enhancements}

While CollabVerse already provides significant value to its users, several enhancements are planned for future releases. The Hopcroft-Karp algorithm could replace Dinic's algorithm for bipartite matching, as it achieves the same asymptotic complexity \(O(E \sqrt{V})\) but with smaller constant factors and simpler implementation. This would improve performance for very large campuses with thousands of students.

The skill ingestion module could be extended beyond language detection to analyze commit frequency, contribution patterns, and collaboration history. This deeper analysis would provide a more nuanced understanding of each student's working style and compatibility with potential teammates. Large Language Models could be employed to generate natural language explanations for why specific students were matched together, helping users understand the rationale behind algorithmic recommendations.

A shared code sandbox integrated directly into the chat interface would enable real-time collaborative coding. This feature could be built using WebRTC for peer-to-peer data transfer and CRDTs (Conflict-free Replicated Data Types) for conflict-free collaborative editing. Finally, organizational support would allow professors and club leads to create verified organizations, run team formation for entire classes, and track project outcomes over multiple semesters.

\newpage
\section{Conclusion}

CollabVerse successfully demonstrates the practical application of advanced algorithms from the Design and Analysis of Algorithms curriculum to a real-world problem of significant practical importance. By implementing Dinic's algorithm for maximum bipartite matching, Bitmask Dynamic Programming for local synergy optimization, and an MLP neural network for skill prediction, the platform solves the team formation problem with provable optimality guarantees. The hybrid architecture, combining FastAPI for web services with a high-performance C++ engine for compute-intensive operations, provides a scalable and maintainable foundation for production deployment.

The project validates several important lessons. Theoretical algorithm knowledge translates directly to impactful software when applied to well-formulated problems. Hybrid architectures can successfully combine the ease of development of high-level languages with the performance characteristics of systems languages like C++. Public APIs, such as GitHub's REST API, provide sufficient data for meaningful inference when cleverly combined with machine learning and algorithmic techniques.

All source code for CollabVerse is available as open source at \url{https://github.com/Cyberpunk-San/CollabVerse}, and the team welcomes contributions from the broader academic community.

\newpage
\section{References}

\begin{enumerate}
    \item Cormen, T. H., Leiserson, C. E., Rivest, R. L., \& Stein, C. (2009). \textit{Introduction to Algorithms} (3rd ed.). MIT Press.
    \item Dinic, E. A. (1970). Algorithm for solution of a problem of maximum flow in a network with power estimation. \textit{Soviet Mathematics Doklady}, 11, 1277–1280.
    \item Even, S., \& Tarjan, R. E. (1975). Network flow and testing graph connectivity. \textit{SIAM Journal on Computing}, 4(4), 507–518.
    \item Hopcroft, J. E., \& Karp, R. M. (1973). An \(n^{5/2}\) algorithm for maximum matchings in bipartite graphs. \textit{SIAM Journal on Computing}, 2(4), 225–231.
    \item Goodfellow, I., Bengio, Y., \& Courville, A. (2016). \textit{Deep Learning}. MIT Press.
    \item GitHub REST API Documentation. \url{https://docs.github.com/en/rest}
    \item FastAPI Documentation. \url{https://fastapi.tiangolo.com/}
    \item Tailwind CSS Documentation. \url{https://tailwindcss.com/docs}
\end{enumerate}

\end{document}